% GENERATED FILE. Edit canonical lessons, then run npm run generate:books.
% canonical-source-hash: fnv1a64:a81a5090f8dcaa70
% canonical-lessons: UR-C15-barish, UR-C15-dhoop, UR-C15-hawa, UR-C15-garmi, UR-C15-sardi

\chapter{Rain, Sun, Wind, Heat, Cold}
\label{ch:ur-weather}
\hlchaptermodality{15}

\begin{chapteropening}
\textbf{By the end of this chapter:} \emph{I can name rain, sunshine, wind, heat and cold in Urdu, and tell which of these words' roots reach back to a documented ancestor and which stop at an honest dead end.}

Say all five weather words in one breath and sort each by how far its root reaches: two Persian dead ends, one inherited dead end, one loan of a loan, and \ur{گرمی}'s own root reaching all the way to Greek thermós — the same two-roads shape \ur{دل} showed on the body's most-named part.
\end{chapteropening}

\section[bārish]{\ur{بارش} — \textquotedblleft{}rain,\textquotedblright{} and why you'd wear the last chapter's coat}

\label{lesson:UR-C15-barish}



\begin{quote}
Recall \textbf{\ur{کوٹ}}, the modern English loan that closed the last chapter. This chapter gives you a reason to wear it: four words for weather, the small talk that follows almost any \textquotedblleft{}how are you.\textquotedblright{}
\end{quote}

\subsection*{You'll want to know first — one word}
\begin{quote}
\textbf{\ur{بارش}} — \emph{bārish} — \textbf{rain}
\end{quote}

\begin{quote}
\textbf{\ur{بارش} \ur{ہو} \ur{رہی} \ur{ہے۔}} — \emph{bārish ho rahī hai.} — \textquotedblleft{}It's raining.\textquotedblright{}
\end{quote}

\textbf{\ur{بارش}} is feminine.

\begin{cousinweb}
\textbf{\ur{بارش}} is borrowed from Classical Persian \textbf{bāriš}, itself built from \textbf{bār}, \textquotedblleft{}load, downpour,\textquotedblright{} plus the ending \textbf{-eš} that turns a verb or noun stem into an action noun. No further Proto-Indo-European root is documented behind that Persian base — an honest dead end, the same shape as \textbf{\ur{لال}}'s.
\end{cousinweb}

\subsection*{Your turn}
\begin{itemize}
\raggedright
  \item \emph{Say it:} \textbf{bārish} — rain; then \textbf{bārish ho rahī hai}
  \item \emph{Name:} \textbf{\ur{بارش}}'s two Persian parts — \textbf{bār} plus the \textbf{-eš} ending
  \item \emph{Connect:} \textbf{\ur{بارش}}'s dead end is the same shape as \textbf{\ur{لال}}'s — Persian, and no further
  \item \emph{Say it:} \textbf{\ur{آج} \ur{بارش} \ur{ہے،} \ur{اس} \ur{لیے} \ur{میرا} \ur{کوٹ} \ur{چاہیے۔}} — \textquotedblleft{}Today there's rain, so I need my coat.\textquotedblright{}
\end{itemize}

\begin{tcolorbox}[breakable,colback=teal!4,colframe=teal!35!black,title={Before you move on}]
What does \textbf{\ur{بارش}} mean, and is it masculine or feminine? (\textbf{Rain}; \textbf{feminine}.) What two Persian pieces build the word? (\textbf{bār, \textquotedblleft{}load/downpour,\textquotedblright{} plus the -eš ending.}) Recall the three respect levels a single Urdu word for \textquotedblleft{}you\textquotedblright{} once had to distinguish — name them. (\textbf{āp, tum, tū.}) Which letter in \textbf{\ur{حافظ}} carries the \emph{f} sound? (\textbf{\ur{ف}.})

Sources: \href{https://en.wiktionary.org/wiki/\%D8\%A8\%D8\%A7\%D8\%B1\%D8\%B4}{Wiktionary: \ur{بارش}}.
\end{tcolorbox}
\section[dhūp]{\ur{دھوپ} — \textquotedblleft{}sunshine,\textquotedblright{} rain's opposite and its etymological twin}

\label{lesson:UR-C15-dhoop}



\begin{quote}
Recall \textbf{\ur{بارش}} — rain, Persian, and a dead end. This word means the opposite in the sky, and reaches a dead end from the other side of the book's two layers.
\end{quote}

\subsection*{You'll want to know first — one word}
\begin{quote}
\textbf{\ur{دھوپ}} — \emph{dhūp} — \textbf{sunshine}
\end{quote}

\begin{quote}
\textbf{\ur{آج} \ur{دھوپ} \ur{ہے۔}} — \emph{āj dhūp hai.} — \textquotedblleft{}There's sunshine today.\textquotedblright{}
\end{quote}

\textbf{\ur{دھوپ}} is feminine.

\begin{cousinweb}
\textbf{\ur{دھوپ}} is inherited from Ashokan Prakrit \textbf{*dhuppā}, \textquotedblleft{}sunshine.\textquotedblright{} No further root is documented behind that. \textbf{\ur{بارش}} hit its dead end on the Persian side; \textbf{\ur{دھوپ}} hits the same kind of honest wall on the inherited side — the same shape \textbf{\ur{کان}} and \textbf{\ur{روٹی}} already showed, not every inherited word traces back to a tidy Indo-European root.
\end{cousinweb}

\subsection*{Your turn}
\begin{itemize}
\raggedright
  \item \emph{Say it:} \textbf{dhūp} — sunshine; then \textbf{āj dhūp hai}
  \item \emph{Contrast:} \textbf{\ur{بارش}}, Persian dead end; \textbf{\ur{دھوپ}}, inherited dead end — opposite weather, same kind of etymological wall
  \item \emph{Name:} two other inherited dead ends already met — \textbf{\ur{کان}} and \textbf{\ur{روٹی}}
  \item \emph{Say it:} \textbf{\ur{آج} \ur{بارش} \ur{ہے} \ur{یا} \ur{دھوپ؟}} — \textquotedblleft{}Is it rain today, or sunshine?\textquotedblright{}
\end{itemize}

\begin{tcolorbox}[breakable,colback=teal!4,colframe=teal!35!black,title={Before you move on}]
What does \textbf{\ur{دھوپ}} mean, and is it masculine or feminine? (\textbf{Sunshine}; \textbf{feminine}.) Is \textbf{\ur{دھوپ}}'s root documented past Prakrit? (\textbf{No — a dead end.}) Recall whether \textbf{\ur{خدا} \ur{حافظ}} is written as one word or two in Urdu spelling. (\textbf{Two, spaced, unlike Persian's joined form.}) What does \textbf{\ur{آنکھ}} mean? (\textbf{Eye.})

Sources: \href{https://en.wiktionary.org/wiki/\%D8\%AF\%DA\%BE\%D9\%88\%D9\%BE}{Wiktionary: \ur{دھوپ}}.
\end{tcolorbox}
\section[hawā]{\ur{ہوا} — \textquotedblleft{}wind,\textquotedblright{} a loan carried twice rather than split once}

\label{lesson:UR-C15-hawa}



\begin{quote}
Recall \textbf{\ur{دھوپ}} — sunshine, an inherited dead end. This word moves with the sunshine and the rain both, and its history takes a shape this book has not shown yet: not two roads from one root, but one word handed across two borrowings in a row.
\end{quote}

\subsection*{You'll want to know first — one word}
\begin{quote}
\textbf{\ur{ہوا}} — \emph{hawā} — \textbf{air, wind}
\end{quote}

\begin{quote}
\textbf{\ur{آج} \ur{ہوا} \ur{چل} \ur{رہی} \ur{ہے۔}} — \emph{āj hawā chal rahī hai.} — \textquotedblleft{}There's wind blowing today.\textquotedblright{}
\end{quote}

\textbf{\ur{ہوا}} is feminine.

\begin{cousinweb}
\textbf{\ur{ہوا}} is borrowed from Classical Persian \emph{hawā} — but Persian did not invent it either. Persian itself borrowed the word from Arabic \textbf{\ur{هَوَاء}} \emph{hawā}. \textbf{\ur{خدا}}, back in the take-leave chapter, came into Urdu directly from Persian, inherited within Iranian the whole way. \textbf{\ur{دل}} and \textbf{\ur{سفید}} each split from one PIE root into two independent roads, Iranian and Indo-Aryan. \textbf{\ur{ہوا}} is a different shape again: a single word passed hand to hand, Arabic to Persian, Persian to Urdu — a loan of a loan, not a shared inheritance at all.
\end{cousinweb}

\subsection*{Your turn}
\begin{itemize}
\raggedright
  \item \emph{Say it:} \textbf{hawā} — wind; then \textbf{āj hawā chal rahī hai}
  \item \emph{Name:} three shapes a word's history can take — one root split two ways (\textbf{\ur{دل}}), a word carried straight through one language (\textbf{\ur{خدا}}), or borrowed twice in a row (\textbf{\ur{ہوا}})
  \item \emph{Say it:} \textbf{\ur{بارش،} \ur{دھوپ،} \ur{اور} \ur{ہوا}} — rain, sunshine, and wind
  \item \emph{Connect:} \textbf{\ur{ہوا}} $\leftarrow$ Persian $\leftarrow$ Arabic \textbf{hawā}, a loan of a loan
\end{itemize}

\begin{tcolorbox}[breakable,colback=teal!4,colframe=teal!35!black,title={Before you move on}]
What does \textbf{\ur{ہوا}} mean, and is it masculine or feminine? (\textbf{Air, wind}; \textbf{feminine}.) How is \textbf{\ur{ہوا}}'s history different from \textbf{\ur{دل}}'s? (\textbf{\ur{دل} split from one root into two roads; \ur{ہوا} is one word borrowed twice in a row.}) Name your sister and your friend in Urdu. (\textbf{\ur{بہن}}; \textbf{\ur{دوست}}.)

Sources: \href{https://en.wiktionary.org/wiki/\%DB\%81\%D9\%88\%D8\%A7}{Wiktionary: \ur{ہوا}}.
\end{tcolorbox}
\section[garmī]{\ur{گرمی} — \textquotedblleft{}heat,\textquotedblright{} a secure cousin and a disputed one}

\label{lesson:UR-C15-garmi}



\begin{quote}
Recall \textbf{\ur{ہوا}} — wind, a loan of a loan. This word is Persian too, but its root behaves like \textbf{\ur{دل}}'s and \textbf{\ur{سفید}}'s: it keeps going back to a shared PIE root, with a genuine cousin waiting on the other side — though this time the cousin needs a careful claim, not a confident one.
\end{quote}

\subsection*{You'll want to know first — one word}
\begin{quote}
\textbf{\ur{گرمی}} — \emph{garmī} — \textbf{heat}
\end{quote}

\begin{quote}
\textbf{\ur{آج} \ur{بہت} \ur{گرمی} \ur{ہے۔}} — \emph{āj bahut garmī hai.} — \textquotedblleft{}It's very hot today.\textquotedblright{}
\end{quote}

\textbf{\ur{گرمی}} is feminine.

\subsection*{The letters in this word}
From the right edge: \textbf{\ur{گ}}, \textbf{\ur{ر}}, \textbf{\ur{م}}, \textbf{\ur{ی}}.

\textbf{\ur{گ}} \emph{gāf} is new: it is \textbf{\ur{ک}} \emph{kāf}, already known, with one extra stroke laid across the top — the same dot-count logic that separated \textbf{\ur{ب}} from \textbf{\ur{پ}}, except this time the letter gains a stroke rather than a dot.

\begin{cousinweb}
\textbf{\ur{گرمی}} continues Middle Persian \emph{garmīh}, from Proto-Iranian \textbf{*garmáh}, from Proto-Indo-Iranian \textbf{*g\textsuperscript{h}armás}, from Proto-Indo-European \textbf{*g\textsuperscript{w}\textsuperscript{h}er-}, \textquotedblleft{}warm.\textquotedblright{} That root securely gives Greek \textbf{thermós}, \textquotedblleft{}hot\textquotedblright{} — the source of English \textbf{thermal} and \textbf{thermometer}. English \textbf{warm} is often mentioned beside it, but dictionaries only call \textbf{warm} a \textquotedblleft{}distant doublet\textquotedblright{} of \emph{thermós}: a proposed relative, not a settled one, with a separate and more likely Germanic root offered as \textbf{warm}'s real source. \textbf{\ur{گرمی}}'s certain family is Greek's, not English's plainest weather word.
\end{cousinweb}

\subsection*{Your turn}
\begin{itemize}
\raggedright
  \item \emph{Say it:} \textbf{garmī} — heat; then \textbf{āj bahut garmī hai}
  \item \emph{Name:} \textbf{\ur{گ}}'s relation to \textbf{\ur{ک}} — one extra stroke, not a dot
  \item \emph{Sort:} secure cousin — Greek \textbf{thermós}; disputed cousin — English \textbf{warm}
  \item \emph{Say it:} \textbf{\ur{بارش،} \ur{دھوپ،} \ur{ہوا،} \ur{گرمی}} — all four weather words so far
\end{itemize}

\begin{tcolorbox}[breakable,colback=teal!4,colframe=teal!35!black,title={Before you move on}]
What does \textbf{\ur{گرمی}} mean, and is it masculine or feminine? (\textbf{Heat}; \textbf{feminine}.) Which English word is \textbf{\ur{گرمی}}'s secure cousin, and which is only disputed? (\textbf{Secure: thermal/thermometer, via Greek thermós. Disputed: warm.}) Name your family in Urdu. (\textbf{\ur{خاندان}.}) What does \textbf{\ur{منہ}} mean? (\textbf{Mouth.})

Sources: \href{https://en.wiktionary.org/wiki/\%DA\%AF\%D8\%B1\%D9\%85\%DB\%8C}{Wiktionary: \ur{گرمی}}, \href{https://en.wiktionary.org/wiki/warm}{Wiktionary: warm}.
\end{tcolorbox}
\section[sardī]{\ur{سردی} — \textquotedblleft{}cold,\textquotedblright{} and one ending that ties this whole book's weather and warmth together}

\label{lesson:UR-C15-sardi}



\begin{quote}
Recall \textbf{\ur{گرمی}} — heat, with a secure cousin in Greek \textbf{thermós} and only a disputed one in English \textbf{warm}. This chapter's last word is its exact opposite, in meaning and in how far its own root travels.
\end{quote}

\subsection*{You'll want to know first — one word}
\begin{quote}
\textbf{\ur{سردی}} — \emph{sardī} — \textbf{cold}
\end{quote}

\begin{quote}
\textbf{\ur{آج} \ur{بہت} \ur{سردی} \ur{ہے۔}} — \emph{āj bahut sardī hai.} — \textquotedblleft{}It's very cold today.\textquotedblright{}
\end{quote}

\textbf{\ur{سردی}} is feminine, exactly like \textbf{\ur{گرمی}}.

\begin{grammarlens}[title={Grammar Lens: one ending, two weather words}]
\textbf{\ur{گرمی}} and \textbf{\ur{سردی}} are built the same way: the plain adjective — \textbf{\ur{گرم}} \emph{garm} \textquotedblleft{}hot,\textquotedblright{} \textbf{\ur{سرد}} \emph{sard} \textquotedblleft{}cold\textquotedblright{} — plus the ending \textbf{\ur{ی}} that turns it into a noun naming the weather itself. Neither word is a one-off; the ending is a small, reusable tool, the same kind of thing \textbf{-\ur{نا}} was for verbs many chapters back.
\end{grammarlens}

\begin{cousinweb}
\textbf{\ur{سردی}} is borrowed from Classical Persian \emph{sardī}, from \emph{sard}, \textquotedblleft{}cold.\textquotedblright{} No further root is documented for that Persian base — a dead end, the same shape as \textbf{\ur{لال}}'s and \textbf{\ur{بارش}}'s, and the mirror image of \textbf{\ur{گرمی}}, whose root reaches all the way to Greek \textbf{thermós}.

Five weather words, every shape this book's etymology has found: \textbf{\ur{بارش}} and \textbf{\ur{سردی}}, Persian dead ends; \textbf{\ur{دھوپ}}, an inherited dead end; \textbf{\ur{ہوا}}, a loan of a loan; \textbf{\ur{گرمی}}, a root reaching a real if disputed cousin — the same two-roads shape \textbf{\ur{دل}} showed on the body's most-named part.
\end{cousinweb}

\subsection*{Your turn}
\begin{itemize}
\raggedright
  \item \emph{Say it:} \textbf{sardī} — cold; then \textbf{āj bahut sardī hai}
  \item \emph{Contrast:} \textbf{\ur{گرمی}}, a root reaching Greek; \textbf{\ur{سردی}}, a Persian dead end
  \item \emph{Say it:} all four weather words in one breath — \textbf{\ur{بارش،} \ur{دھوپ،} \ur{ہوا،} \ur{گرمی،} \ur{سردی}}
  \item \emph{Recall:} \textbf{\ur{دل}}'s two roads from one PIE root, and how \textbf{\ur{گرمی}}'s own root follows the same shape
\end{itemize}

\begin{tcolorbox}[breakable,colback=teal!4,colframe=teal!35!black,title={Before you move on}]
What does \textbf{\ur{سردی}} mean? (\textbf{Cold}.) Does its root reach as far as \textbf{\ur{گرمی}}'s? (\textbf{No — a Persian dead end.}) Which letter in \textbf{\ur{بھائی}} marks the glide between two vowels? (\textbf{\ur{ئ}, ye with hamza.}) Which two face words share the same three letters reversed? (\textbf{\ur{کان} and \ur{ناک}.}) What job does plain \textbf{\ur{ن}} do in \textbf{\ur{آنکھ}} and \textbf{\ur{منہ}}? (\textbf{Marks the vowel before it as nasal.}) What PIE root ties \textbf{\ur{دل}} to Sanskrit \textbf{hṛdaya}? (\textbf{*\'{k}ērd-}, split at the Indo-Iranian fork.)

Sources: \href{https://en.wiktionary.org/wiki/\%D8\%B3\%D8\%B1\%D8\%AF\%DB\%8C}{Wiktionary: \ur{سردی}}.
\end{tcolorbox}
